\section{Small-Budget Theory for Non-Adaptive Payoffs}
\label{sec:small-budget-first-stage-objectives-appendix}

This section treats the small-budget problem, in \Cref{sec:generic-first-stage-results-appendix} for non-adaptive first-stage objectives of the form $V_Y(\bm s_1)=\E[\Test_1\cdot\, Y]$, and in \Cref{sec:named-first-stage-results-appendix} for the named objectives from the paper by specialization.

\subsection{Generic first-stage results}
\label{sec:generic-first-stage-results-appendix}

The following gives the proof of the generic small-investment expansion.
\refstepcounter{proposition}
\begin{linkedresult}{proof}{prop:small-investment-expansion}
  \label{prop:small-investment-expansion-appendix}
  Fix a prior distribution of $\bm\tau$, and let $Y$ be non-adaptive with $\E[|Y|]<\infty$ and $\E[|Y|\|\bm\tau\|_2^3]<\infty$.
  Then there exist constants $\overline{\invst}>0$ and $0<C_0,C_1<\infty$ such that for all pilot investments $\invst\in(0,\overline{\invst}]$ and all unit-investment samples $\overline{\bm{s}}_1\in\overline{\mathcal{S}}$,
  \[
    V_Y(\invst\,\overline{\bm{s}}_1)
    =
    \alpha_1 \E[Y]
    +\sqrt{\invst}\,\varphi(z_{1-\alpha_1})\,
    \frac{(\bm{\theta}^Y)^{\!\top}\overline{\bm{s}}_1}{\sqrt{\bm{1}^{\top}\overline{\bm{s}}_1}}
    +R(\invst,\overline{\bm{s}}_1).
  \]

  The remainder satisfies the following for all $\invst, \invst'\in(0,\overline{\invst}]$ and all $\overline{\bm{s}}_1,\overline{\bm{s}}_1'\in\overline{\mathcal{S}}$.
  It is bounded as
  \[
    \left|R(\invst,\overline{\bm{s}}_1)\right|
    \le C_0\,\invst,
  \]
  it is Lipschitz in the investment $\invst$
  \[
    \left|R(\invst,\overline{\bm{s}}_1)-R(\invst',\overline{\bm{s}}_1)\right|
    \le C_0\,|\invst-\invst'|,
  \]
  and it is $O(\invst)$-Lipschitz in the unit-investment design
  \[
    \left|R(\invst,\overline{\bm{s}}_1)-R(\invst,\overline{\bm{s}}_1')\right|
    \le C_1\,\invst\,\|\overline{\bm{s}}_1-\overline{\bm{s}}_1'\|_1.
  \]
\end{linkedresult}

\begin{proof}
  First, we give a smooth extension and reparametrization of the first-stage objective around $\invst=0$.
  Then we use a Taylor expansion to calculate the first two terms.
  Finally, we bound the remainder term and show it has the desired properties.

  \emph{Step 1: A $C^{\infty}$ reparametrization around $\invst=0$.}
  Reparametrize the pilot sample $\bm{s}_1 = j^2 \overline{\bm{s}}_1$ with $j\in\R$ and $\overline{\bm{s}}_1\in \overline{\overline{\mathcal{S}}^{\epsilon}}$ where $\overline{\overline{\mathcal{S}}^{\epsilon}}\supset \overline{\mathcal{S}}^{\epsilon}$ are closed and open $\epsilon$-neighborhoods respectively of the set of unit-investment samples $\overline{\mathcal{S}}$ within $\R_{\ge 0}^D$, with $\epsilon = 1 / (2 \norm{\bm{c}}_2) > 0$.
  For $\overline{\bm{s}}_{1}\in\overline{\overline{\mathcal{S}}^{\epsilon}}$ define
  \[
    \overline{n}_{1} \defeq \bm{1}^{\top}\overline{\bm{s}}_{1}.
  \]
  Note that for all $\overline{\bm{s}}_{1}\in\overline{\overline{\mathcal{S}}^{\epsilon}}$ we can write $\overline{\bm{s}}_1 = \overline{\bm{s}}_1' + \bm{\epsilon}$ for some $\overline{\bm{s}}_1'\in\overline{\mathcal{S}}$ and $\bm{\epsilon}$ with $\norm{\bm{\epsilon}}_2 \le \epsilon$.
  Therefore
  \[
    \bm{c}^\top \overline{\bm{s}}_1 = 1 + \bm{c}^\top \bm{\epsilon}
    \ge 1 - \norm{\bm{c}}_2 \epsilon = 1/2.
  \]
  Since $\overline{\bm{s}}_1\in\R_{\ge 0}^D$, Hölder's inequality gives
  \[
    \left|\bm{c}^\top \overline{\bm{s}}_1\right| \le
    \norm{\bm{c}}_\infty
    \norm{\overline{\bm{s}}_1}_1,
    \qquad
    \left|\bm{c}^\top \overline{\bm{s}}_1\right| \le \norm{\bm{c}}_2
    \norm{\overline{\bm{s}}_1}_2.
  \]
  Combining these, we have
  \[
    1/4 \le \norm{\bm{c}}_2^2 \norm{\overline{\bm{s}}_1}_2^2,
    \qquad
    1/2 \le \norm{\bm{c}}_\infty \norm{\overline{\bm{s}}_1}_1.
  \]
  Thus
  \[
    \overline{n}_{1}
    =
    \norm{\overline{\bm{s}}_{1}}_1
    \ge
    \frac{1}{2 \norm{\bm{c}}_\infty}
    > 0.
  \]

  Define $\xi:\overline{\overline{\mathcal{S}}^{\epsilon}}\times\R^D\to\R$ and $q:\R\times\overline{\overline{\mathcal{S}}^{\epsilon}}\times\R^D\to[0,1]$ as
  \[
    \xi(\overline{\bm{s}}_1,\bm\tau)
    \defeq
    \frac{\overline{\bm{s}}_1^\top\bm\tau}{\sqrt{\overline n_1}},
    \qquad
    q(j,\overline{\bm{s}}_1,\bm\tau)
    \defeq
    \Phi\big(j\,\xi(\overline{\bm{s}}_1,\bm\tau)-z_{1-\alpha_1}\big),
  \]
  such that if $j>0$ and $\bm s_1=j^2\overline{\bm s}_1$, then
  \[
    q(j,\overline{\bm{s}}_1,\bm\tau)
    =
    \P(\Test_1=1\mid \bm\tau).
  \]
  Since $Y$ is assumed to be non-adaptive---that is, independent of the pilot-stage sampling noise conditional on $\bm\tau$---iterated expectations gives
  \begin{align*}
    V_Y(j^2\overline{\bm s}_1)
    & =
    \E\!\left[\Test_1 \cdot \, Y\right] \\
    & =
    \E\!\left[\E\!\left[\Test_1 \mid \bm\tau\right]\cdot \, Y\right] \\
    & =
    \E\!\left[q(j,\overline{\bm{s}}_1,\bm\tau)\,Y\right].
  \end{align*}
  Define the composite map $g:\R\times\overline{\overline{\mathcal{S}}^{\epsilon}}\to\R$ as
  \[
    g(j,\overline{\bm s}_1)
    \defeq
    \E\!\left[q(j,\overline{\bm{s}}_1,\bm\tau)\,Y\right]
  \]
  so that $V_Y(j^2\overline{\bm s}_1)=g(j,\overline{\bm s}_1)$ for $j>0$.

  Since $\overline{\overline{\mathcal{S}}^\epsilon}$ is compact, there exists $M<\infty$ such that
  \[
    \|\overline{\bm s}_1\|_2\le M
    \qquad\forall\overline{\bm s}_1\in \overline{\overline{\mathcal{S}}^\epsilon}.
  \]
  Therefore
  \[
    |\xi(\overline{\bm{s}}_1,\bm\tau)|
    \le
    \frac{M}{\sqrt{\overline n_1}}\|\bm\tau\|_2
    \le
    M\sqrt{2\|\bm c\|_\infty}\,\|\bm\tau\|_2.
  \]
  Also, differentiating $\xi$ with respect to $\overline{\bm s}_1$ gives
  \[
    \nabla_{\overline{\bm s}_1}\xi(\overline{\bm s}_1,\bm\tau)
    =
    \frac{\bm\tau}{\sqrt{\bm 1^\top\overline{\bm s}_1}}
    -
    \frac{\overline{\bm s}_1^\top\bm\tau}{2(\bm 1^\top\overline{\bm s}_1)^{3/2}}\,\bm 1,
  \]
  and hence, using $\|\bm 1\|_\infty=1$,
  \[
    \left\|\nabla_{\overline{\bm s}_1}\xi(\overline{\bm s}_1,\bm\tau)\right\|_\infty
    \le
    \left(\sqrt{2\|\bm c\|_\infty}+\sqrt{2}M\|\bm c\|_\infty^{3/2}\right)\|\bm\tau\|_2.
  \]

  Fix arbitrary $\overline j>0$, and write $\overline{\overline{\mathcal{S}}^{\epsilon/2}}$ for the closed $\epsilon/2$-neighborhood of $\overline{\mathcal{S}}$.
  The first two derivatives of $q$ with respect to $j$ are
  \[
    \frac{\partial}{\partial j}q(j,\overline{\bm s}_1,\bm\tau)
    =
    \xi(\overline{\bm s}_1,\bm\tau)\,
    \varphi\!\big(j\,\xi(\overline{\bm s}_1,\bm\tau) - z_{1-\alpha_1}\big),
  \]
  and
  \[
    \frac{\partial^2}{\partial j^2}q(j,\overline{\bm s}_1,\bm\tau)
    =
    - \xi(\overline{\bm s}_1,\bm\tau)^2
    \big(j\,\xi(\overline{\bm s}_1,\bm\tau) - z_{1-\alpha_1}\big)
    \varphi\!\big(j\,\xi(\overline{\bm s}_1,\bm\tau) - z_{1-\alpha_1}\big).
  \]
  Since $\varphi$, $x\varphi(x)$, and $x^2\varphi(x)$ are bounded on $\R$ ($\varphi$ is bounded and its tails decay exponentially, i.e. faster than any polynomial), there exists a constant $C<\infty$ depending only on $\overline j$, $\bm c$, $M$, and $z_{1-\alpha_1}$ such that for all $(j,\overline{\bm s}_1,\bm\tau)\in[-\overline j,\overline j]\times \overline{\overline{\mathcal{S}}^{\epsilon}}\times\R^D$,
  \[
    \left|\frac{\partial}{\partial j}q(j,\overline{\bm s}_1,\bm\tau)\right|
    \le
    C\,\|\bm\tau\|_2,
    \qquad
    \left|\frac{\partial^2}{\partial j^2}q(j,\overline{\bm s}_1,\bm\tau)\right|
    \le
    C\,\|\bm\tau\|_2^2.
  \]
  Writing
  \[
    h_j(x)\defeq - x^2(jx - z_{1-\alpha_1})\varphi(jx - z_{1-\alpha_1}),
  \]
  we have
  \[
    \frac{\partial^2}{\partial j^2}q(j,\overline{\bm s}_1,\bm\tau)
    =
    h_j\!\big(\xi(\overline{\bm s}_1,\bm\tau)\big),
  \]
  and the chain rule gives
  \[
    \nabla_{\overline{\bm s}_1}\frac{\partial^2}{\partial j^2}q(j,\overline{\bm s}_1,\bm\tau)
    =
    h_j'\!\big(\xi(\overline{\bm s}_1,\bm\tau)\big)\,
    \nabla_{\overline{\bm s}_1}\xi(\overline{\bm s}_1,\bm\tau).
  \]
  Since
  \[
    h_j'(x)
    =
    -2x(jx-z_{1-\alpha_1})\varphi(jx-z_{1-\alpha_1})
    -jx^2\varphi(jx-z_{1-\alpha_1})
    +jx^2(jx-z_{1-\alpha_1})^2\varphi(jx-z_{1-\alpha_1}),
  \]
  and $\varphi$, $x\varphi(x)$, and $x^2\varphi(x)$ are bounded, there exists $C'<\infty$ such that
  \[
    |h_j'(x)|\le C'(|x|+x^2)
    \qquad
    \forall (j,x)\in[-\overline j,\overline j]\times\R.
  \]
  Combining this with the bounds on $\xi$ and $\nabla_{\overline{\bm s}_1}\xi$ yields
  \[
    \left\|\nabla_{\overline{\bm s}_1}\frac{\partial^2}{\partial j^2}q(j,\overline{\bm s}_1,\bm\tau)\right\|_\infty
    \le
    C''\left(\|\bm\tau\|_2^2+\|\bm\tau\|_2^3\right)
  \]
  for a constant $C''<\infty$.

  By the integrability conditions assumed on $Y$ and $\bm{\tau}$,
  \[
    2\left(\mathbb{E}[|Y|]+\mathbb{E}\!\left[|Y|\|\boldsymbol{\tau}\|_2^3\right]\right)
    =
    \E\left[2|Y|\left(1+\|\boldsymbol{\tau}\|_2^3\right)\right]
    <
    \infty.
  \]
  Also, for \(x\ge 0\),
  \[
    x,\ x^2,\ x^2+x^3 \le 2(1+x^3).
  \]
  Thus, applying this with \(x=\|\boldsymbol{\tau}\|_2\) and multiplying by \(|Y|\), the random variables
  \[
    |Y|\,\|\boldsymbol{\tau}\|_2,
    \qquad
    |Y|\,\|\boldsymbol{\tau}\|_2^2,
    \qquad
    |Y|\left(\|\boldsymbol{\tau}\|_2^2+\|\boldsymbol{\tau}\|_2^3\right)
  \]
  are integrable.

  Moreover, for each fixed $\bm\tau$, the maps
  \[
    (j,\overline{\bm s}_1)\mapsto q(j,\overline{\bm s}_1,\bm\tau),
    \quad
    \frac{\partial}{\partial j}q(j,\overline{\bm s}_1,\bm\tau),
    \quad
    \frac{\partial^2}{\partial j^2}q(j,\overline{\bm s}_1,\bm\tau),
    \quad
    \nabla_{\overline{\bm s}_1}\frac{\partial^2}{\partial j^2}q(j,\overline{\bm s}_1,\bm\tau)
  \]
  are continuous on $[-\overline j,\overline j]\times \overline{\overline{\mathcal{S}}^\epsilon}$.
  Therefore, we can differentiate
  \[
    g(j,\overline{\bm{s}}_1) = \E[q(j,\overline{\bm{s}}_1,\bm{\tau}) Y]
  \]
  under the expectation~\cite[Theorem~2.27]{Folland99}, first in $j$ and then componentwise in $\overline{\bm s}_1$,
  \[
    \frac{\partial}{\partial j}g(j,\overline{\bm s}_1)
    =
    \E\!\left[
      Y\,
      \frac{\partial}{\partial j}q(j,\overline{\bm s}_1,\bm\tau)
    \right],
    \qquad
    \frac{\partial^2}{\partial j^2}g(j,\overline{\bm s}_1)
    =
    \E\!\left[
      Y\,
      \frac{\partial^2}{\partial j^2}q(j,\overline{\bm s}_1,\bm\tau)
    \right]
  \]
  for $(j,\overline{\bm s}_1)\in(-\overline j,\overline j)\times \overline{\mathcal{S}}^\epsilon$, and
  \[
    \nabla_{\overline{\bm s}_1}\frac{\partial^2}{\partial j^2}g(j,\overline{\bm s}_1)
    =
    \E\!\left[
      Y\,
      \nabla_{\overline{\bm s}_1}\frac{\partial^2}{\partial j^2}q(j,\overline{\bm s}_1,\bm\tau)
    \right]
  \]
  on the same set.
  All of these derivatives are continuous by dominated convergence: the preceding bounds give parameter-uniform integrable envelopes for the integrands, while the lower bound on \(\bm 1^\top\overline{\bm s}_1\) ensures pointwise continuity in \((j,\overline{\bm s}_1)\).

  Finally, since
  \[
    \mathcal D\defeq[0,\overline j / 2]\times \overline{\overline{\mathcal{S}}^{\epsilon/2}}
  \]
  is a compact subset of $(-\overline j,\overline j)\times\overline{\mathcal{S}}^{\epsilon}$, we may define
  \[
    C_0\defeq\frac{1}{2}\sup_{(j,\overline{\bm s}_1)\in\mathcal D}\left|
    \frac{\partial^2}{\partial j^2}g(j,\overline{\bm s}_1)
    \right|<\infty,
    \qquad
    C_1\defeq\frac{1}{2}\sup_{(j,\overline{\bm s}_1)\in\mathcal D}\left\|
    \nabla_{\overline{\bm s}_1}\frac{\partial^2}{\partial j^2}g(j,\overline{\bm s}_1)
    \right\|_\infty<\infty.
  \]

  \emph{Step 2: Constant term and first derivative at $j=0$.}
  At $j=0$, for any $\overline{\bm s}_1\in\overline{\mathcal{S}}$, we have $q(0,\overline{\bm s}_1,\bm\tau)=1 - \Phi(z_{1-\alpha_1})=\alpha_1$, so
  \[
    g(0,\overline{\bm s}_1)=\alpha_1\,\E[Y].
  \]
  Evaluating the formula above for $\partial g/\partial j$ at $j=0$ and plugging in the definition of $\bm{\theta}^Y=\E[\bm\tau Y]$ gives
  \begin{align*}
    \frac{\partial}{\partial j}g(0,\overline{\bm s}_1)
    &=
    \E\!\left[
      \frac{\partial}{\partial j}q(j,\overline{\bm s}_1,\bm\tau)\Big|_{j=0}\,Y
    \right] \\
    &=
    \varphi(z_{1-\alpha_1})\,
    \frac{1}{\sqrt{\bm 1^\top\overline{\bm s}_1}}\,
    \E\!\left[
      (\overline{\bm s}_1^\top\bm\tau)\,Y
    \right] \\
    &=
    \varphi(z_{1-\alpha_1})\,
    \frac{(\bm{\theta}^Y)^\top\overline{\bm s}_1}{\sqrt{\bm 1^\top\overline{\bm s}_1}}.
  \end{align*}

  \emph{Step 3: Taylor expansion with an integral remainder and remainder bounds.}
  For any $(j,\overline{\bm s}_1)\in\mathcal D$, Taylor's theorem with an integral remainder~\cite[Theorem~2.55]{Folland23} gives
  \[
    g(j,\overline{\bm{s}}_1)
    =
    g(0,\overline{\bm{s}}_1)
    +j\,\frac{\partial}{\partial j}g(0,\overline{\bm{s}}_1)
    +\int_0^j (j-x)\,\frac{\partial^2}{\partial j^2}
    g(x,\overline{\bm{s}}_1)\,dx.
  \]
  Let $\overline{\invst}\defeq (\overline{j} / 2)^2$.
  Plugging in the values computed above, for any $\invst\in(0,\overline{\invst}]$ and $\overline{\bm s}_1\in\overline{\mathcal{S}}$,
  \[
    V_Y(\invst\overline{\bm{s}}_1)
    =
    g(\sqrt{\invst},\overline{\bm{s}}_1)
    =
    \alpha_1 \E[Y]
    +\sqrt{\invst}\,\varphi(z_{1-\alpha_1})\,
    \frac{(\bm{\theta}^Y)^{\!\top}\overline{\bm{s}}_1}{\sqrt{\bm{1}^{\top}\overline{\bm{s}}_1}}
    +R(\invst,\overline{\bm{s}}_1),
  \]
  where
  \[
    R(\invst,\overline{\bm{s}}_1)
    \defeq
    \int_0^{\sqrt{\invst}} (\sqrt{\invst}-x)\,\frac{\partial^2}{\partial j^2}
    g(x,\overline{\bm{s}}_1)\,dx.
  \]

  Since $|\frac{\partial^2}{\partial j^2} g|\le 2 C_0$ on $\mathcal D$, the remainder is bounded as
  \[
    |R(\invst,\overline{\bm{s}}_1)|
    \le
    2 C_0 \int_0^{\sqrt{\invst}}(\sqrt{\invst}-x)\,dx
    =
    C_0\,\invst.
  \]

  To show the Lipschitz bound in the investment, fix $\overline{\bm{s}}_1\in\overline{\overline{\mathcal{S}}^{\epsilon/2}}$ and $\invst,\invst'\in(0,\overline{\invst}]$.
  Without loss of generality assume $\invst'\le \invst$.
  Subtracting and splitting the integral at $\sqrt{\invst'}$ gives
  \begin{align*}
    R(\invst,\overline{\bm{s}}_1)-R(\invst',\overline{\bm{s}}_1)
    &=
    (\sqrt{\invst}-\sqrt{\invst'})\int_0^{\sqrt{\invst'}}\frac{\partial^2}{\partial
    j^2}g(x,\overline{\bm{s}}_1)\,dx
    \;+\;
    \int_{\sqrt{\invst'}}^{\sqrt{\invst}}(\sqrt{\invst}-x)\,\frac{\partial^2}{\partial
    j^2}g(x,\overline{\bm{s}}_1)\,dx.
  \end{align*}
  Since $\big|\frac{\partial^2}{\partial j^2}g(x,\overline{\bm{s}}_1)\big|\le 2C_0$ uniformly on $\mathcal{D}$, we can bound
  \begin{align*}
    \big|R(\invst,\overline{\bm{s}}_1)-R(\invst',\overline{\bm{s}}_1)\big|
    &\le
    2C_0\Big((\sqrt{\invst}-\sqrt{\invst'})\sqrt{\invst'}+\int_{\sqrt{\invst'}}^{\sqrt{\invst}}(\sqrt{\invst}-x)\,dx\Big)
    \\
    &=
    2C_0\Big((\sqrt{\invst}-\sqrt{\invst'})\sqrt{\invst'}+\tfrac{1}{2}(\sqrt{\invst}-\sqrt{\invst'})^2\Big)
    \\
    &=
    C_0\,(\invst-\invst').
  \end{align*}
  Therefore, for all $\invst,\invst'\in(0,\overline{\invst}]$ and all $\overline{\bm{s}}_1\in\overline{\overline{\mathcal{S}}^{\epsilon/2}}$,
  \[
    \big|R(\invst,\overline{\bm{s}}_1)-R(\invst',\overline{\bm{s}}_1)\big|
    \le C_0\,|\invst-\invst'|.
  \]

  To show the Lipschitz bound in the unit-investment design, fix $\overline{\bm{s}}_1,\overline{\bm{s}}_1'\in\overline{\overline{\mathcal{S}}^{\epsilon/2}}$ and $x\in(0,\overline{j} / 2]$.
  Note that $\overline{\overline{\mathcal{S}}^{\epsilon/2}}$ is convex because it is the closed $\epsilon/2$-neighborhood of the convex polytope $\overline{\mathcal{S}}$ within $\R_{\ge 0}^D$.
  Therefore, since $\frac{\partial^2}{\partial j^2} g(x,\cdot)$ is continuously differentiable on $\overline{\mathcal{S}}^{\epsilon}\supset \overline{\overline{\mathcal{S}}^{\epsilon/2}}$, the fundamental theorem of calculus along the line segment gives
  \[
    \frac{\partial^2}{\partial j^2}
    g(x,\overline{\bm{s}}_1)-\frac{\partial^2}{\partial j^2}
    g(x,\overline{\bm{s}}_1')
    =
    \int_0^1
    \nabla_{\overline{\bm{s}}_1}\frac{\partial^2}{\partial j^2}
    g\big(x,\overline{\bm{s}}_1'+\lambda(\overline{\bm{s}}_1-\overline{\bm{s}}_1')\big)^{\!\top}
    (\overline{\bm{s}}_1-\overline{\bm{s}}_1')\,d\lambda.
  \]
  Thus, by Hölder's inequality and the bound $\|\nabla_{\overline{\bm{s}}_1}\frac{\partial^2}{\partial j^2} g\|_\infty\le 2 C_1$ on $\mathcal D$, we can bound
  \[
    \left|\frac{\partial^2}{\partial j^2}
    g(x,\overline{\bm{s}}_1)-\frac{\partial^2}{\partial j^2}
    g(x,\overline{\bm{s}}_1')\right|
    \le
    2 C_1\,\|\overline{\bm{s}}_1-\overline{\bm{s}}_1'\|_1.
  \]
  Therefore, for $\invst \le \overline{\invst}$,
  \begin{align*}
    |R(\invst,\overline{\bm{s}}_1)-R(\invst,\overline{\bm{s}}_1')|
    &\le
    \int_0^{\sqrt{\invst}} (\sqrt{\invst}-x)\,
    \left|\frac{\partial^2}{\partial j^2}
    g(x,\overline{\bm{s}}_1)-\frac{\partial^2}{\partial j^2}
    g(x,\overline{\bm{s}}_1')\right|\,dx\\
    &\le
    2 C_1\,\|\overline{\bm{s}}_1-\overline{\bm{s}}_1'\|_1
    \int_0^{\sqrt{\invst}} (\sqrt{\invst}-x)\,dx
    \;=\;
    C_1\,\invst\,\|\overline{\bm{s}}_1-\overline{\bm{s}}_1'\|_1.
  \end{align*}
\end{proof}

The following gives the proof of the generic small-budget optimal design result.

\refstepcounter{theorem}
\begin{linkedresult}{proof}{thm:optimal-small-budget-design-general}
  \label{thm:optimal-small-budget-design-general-appendix}
  Fix a prior distribution of $\bm\tau$ and a payoff $Y$ satisfying the assumptions of
  \linkedCref{prop:small-investment-expansion-appendix}.
  Suppose there exists a unique type $d^\star$ maximizing the small-budget index
  \[
    \pi_d^Y
    \defeq
    \frac{\theta_d^Y}{\sqrt{c_{d}}},
  \]
  and that $\pi_{d^\star}^Y>0$.
  Then there exists $\overline{b}>0$ such that for every budget $b\in(0,\overline{b}]$, a unique optimal design exists for
  \[
    \max_{\substack{\bm s_1\in\R_{\geq 0}^{D} \setminus \{\bm{0}\}\\ \bm c^\top \bm s_1\le b}}
    V_Y(\bm s_1),
  \]
  and it invests the entire budget in the type $d^\star$,
  \[
    \bm{s}_1^\star(b) =
    \frac{b}{c_{d^\star}}\,\bm{e}_{d^\star}.
  \]
\end{linkedresult}

\begin{proof}
  \emph{Basic definitions and outline.}
  Recall from \Cref{sec:Model} that the budget-constrained problem is equivalent to maximizing $V_Y(\invst\,\overline{\bm{s}}_1)$ over $\invst\in(0,b]$ and $\overline{\bm{s}}_1\in\overline{\mathcal{S}}$.
  We proceed in three steps.
  First, we show that the leading-order function
  \[
    L(\overline{\bm{s}}_1)\defeq
    \frac{(\bm{\theta}^Y)^\top
    \overline{\bm{s}}_1}{\sqrt{\bm{1}^\top\overline{\bm{s}}_1}},
    \qquad \overline{\bm{s}}_1\in\overline{\mathcal{S}},
  \]
  in \linkedCref{prop:small-investment-expansion-appendix} is uniquely maximized over $\overline{\mathcal{S}}$ at the vertex corresponding to $d^\star$.
  Second, we show that for sufficiently small fixed investment $\invst$, the unique optimal unit-investment design is the same vertex.
  Third, we show that the optimal design for sufficiently small budgets $b$ invests the entire budget there.

  \emph{Step 1: The leading-order functional $L$ is uniquely maximized over $\overline{\mathcal{S}}$ at the vertex corresponding to $d^\star$ and decreases at least linearly away from it.}
  Define the corresponding unit-investment vertex
  \[
    \overline{\bm{v}}^\star \defeq
    \frac{1}{c_{d^\star}}\,\bm{e}_{d^\star}\in\overline{\mathcal{S}}.
  \]
  Then $L(\overline{\bm{v}}^\star)=\pi_{d^\star}^Y$.
  Define the index gap
  \[
    \Delta\defeq \pi_{d^\star}^Y-\max_{k\neq d^\star}\pi_k^Y \;>\;0,
  \]
  and the cost extrema $c_{\min}\defeq\min_d c_{d}$ and $c_{\max}\defeq\max_d c_{d}$.
  For any $\overline{\bm{s}}_1\in\overline{\mathcal{S}}$ define $\lambda_d\defeq c_{d}\overline{s}_{1d}$, so $\lambda_d\ge0$ and $\sum_d\lambda_d=1$.
  Also, write $w_d\defeq 1/c_{d}$, so that $\overline{s}_{1d}=\lambda_d w_d$.

  Under this new parametrization, we have
  \[
    \bm{1}^\top\overline{\bm{s}}_1
    =
    \sum_{d=1}^D \overline{s}_{1d}
    =
    \sum_{d=1}^D \lambda_d w_d,
    \qquad
    (\bm{\theta}^Y)^\top\overline{\bm{s}}_1
    =
    \sum_{d=1}^D \theta_d^Y\,\overline{s}_{1d}
    =
    \sum_{d=1}^D \lambda_d\,\theta_d^Y\,w_d.
  \]
  Using $\theta_d^Y = \pi_d^Y\sqrt{c_{d}}$ and $w_d=1/c_{d}$ gives $\theta_d^Y w_d = \pi_d^Y\sqrt{w_d}$, hence
  \[
    (\bm{\theta}^Y)^\top\overline{\bm{s}}_1
    =
    \sum_{d=1}^D \lambda_d\,\pi_d^Y\,\sqrt{w_d},
    \qquad\text{and therefore}\qquad
    L(\overline{\bm{s}}_1)
    =
    \frac{\sum_{d=1}^D
    \lambda_d\,\pi_d^Y\,\sqrt{w_d}}{\sqrt{\sum_{d=1}^D \lambda_d w_d}}.
  \]

  By the definition of $\Delta$, for every $d\neq d^\star$ we have $\pi_d^Y\le \pi_{d^\star}^Y-\Delta$, so
  \begin{align*}
    \sum_{d=1}^D \lambda_d\,\pi_d^Y\,\sqrt{w_d}
    &=
    \lambda_{d^\star} \pi_{d^\star}^Y\sqrt{w_{d^\star}}
    +\sum_{d\neq d^\star}\lambda_d\,\pi_d^Y\,\sqrt{w_d} \\
    &\le
    \lambda_{d^\star} \pi_{d^\star}^Y\sqrt{w_{d^\star}}
    +\sum_{d\neq d^\star}\lambda_d\,(\pi_{d^\star}^Y-\Delta)\,\sqrt{w_d} \\
    &=
    \pi_{d^\star}^Y\sum_{d=1}^D \lambda_d\sqrt{w_d}
    -\Delta\sum_{d\neq d^\star}\lambda_d\sqrt{w_d}.
  \end{align*}
  Dividing by $\sqrt{\sum_d \lambda_d w_d}$ yields
  \[
    L(\overline{\bm{s}}_1)
    \le
    \pi_{d^\star}^Y\,
    \frac{\sum_{d=1}^D \lambda_d\sqrt{w_d}}{\sqrt{\sum_{d=1}^D \lambda_d w_d}}
    \;-\;
    \Delta\,
    \frac{\sum_{d\neq
    d^\star}\lambda_d\sqrt{w_d}}{\sqrt{\sum_{d=1}^D \lambda_d w_d}}.
  \]
  By Cauchy--Schwarz,
  \[
    \sum_{d=1}^D \lambda_d\sqrt{w_d}
    =
    \sum_{d=1}^D \sqrt{\lambda_d}\,\sqrt{\lambda_d w_d}
    \le
    \sqrt{\sum_{d=1}^D \lambda_d}\,\sqrt{\sum_{d=1}^D \lambda_d w_d}
    =
    \sqrt{\sum_{d=1}^D \lambda_d w_d},
  \]
  hence the ratio in the first term is at most $1$, and since $\pi_{d^\star}^Y>0$ we obtain
  \[
    L(\overline{\bm{s}}_1)
    \le
    \pi_{d^\star}^Y
    \;-\;
    \Delta\,
    \frac{\sum_{d\neq
    d^\star}\lambda_d\sqrt{w_d}}{\sqrt{\sum_{d=1}^D \lambda_d w_d}}.
  \]
  Noting that $L(\overline{\bm{v}}^\star)=\pi_{d^\star}^Y$, we can rearrange this as
  \[
    L(\overline{\bm{v}}^\star)-L(\overline{\bm{s}}_1)
    \ge
    \Delta\,
    \frac{\sum_{d\neq
    d^\star}\lambda_d\sqrt{w_d}}{\sqrt{\sum_{d=1}^D \lambda_d w_d}}.
  \]
  Since $\min_d \sqrt{w_d}=1/\sqrt{c_{\max}}$ and $\max_d w_d = 1/c_{\min}$,
  \[
    \sum_{d\neq d^\star}\lambda_d\sqrt{w_d}
    \ge
    \frac{1-\lambda_{d^\star}}{\sqrt{c_{\max}}},
    \qquad
    \sqrt{\sum_{d=1}^D \lambda_d w_d}
    \le
    \frac{1}{\sqrt{c_{\min}}}.
  \]
  Combining these gives the uniform bound
  \[
    L(\overline{\bm{v}}^\star)-L(\overline{\bm{s}}_1)
    \ge
    \Delta\sqrt{\frac{c_{\min}}{c_{\max}}}\,(1-\lambda_{d^\star}),
    \qquad \forall\overline{\bm{s}}_1\in\overline{\mathcal{S}}.
  \]
  Finally, we can relate $1-\lambda_{d^\star}$ to the $\ell_1$-distance from the vertex:
  \begin{align*}
    \norm{\overline{\bm{s}}_1-\overline{\bm{v}}^\star}_1
    &=
    \left|\overline{s}_{1d^\star}-\overline{v}_{1d^\star}^\star\right|
    +\sum_{d\neq d^\star}\left|\overline{s}_{1d}-\overline{v}_{1d}^\star\right| \\
    &=
    \left|\overline{s}_{1d^\star}-\frac{1}{c_{d^\star}}\right|
    +\sum_{d\neq d^\star}\left|\overline{s}_{1d}-0\right|
    && \text{since }\overline{\bm{v}}^\star=\bm{e}_{d^\star}/c_{d^\star} \\
    &=
    \frac{1}{c_{d^\star}}-\overline{s}_{1d^\star}
    +\sum_{d\neq d^\star}\overline{s}_{1d}
    && \text{since }c_{d^\star}\overline{s}_{1d^\star}=\lambda_{d^\star}\le 1 \\
    &=
    \frac{1}{c_{d^\star}}-\frac{\lambda_{d^\star}}{c_{d^\star}}
    +\sum_{d\neq d^\star}\frac{\lambda_d}{c_d} \\
    &\le
    \frac{1-\lambda_{d^\star}}{c_{\min}}
    +\sum_{d\neq d^\star}\frac{\lambda_d}{c_{\min}} \\
    &=
    \frac{2}{c_{\min}}\,(1-\lambda_{d^\star}).
  \end{align*}
  Combining this with the previous inequality gives
  \begin{equation}
    \label{eq:generic-leading-gap}
    L(\overline{\bm{v}}^\star)-L(\overline{\bm{s}}_1)
    \ge
    q\,\|\overline{\bm{s}}_1-\overline{\bm{v}}^\star\|_1,
    \qquad
    q\defeq \frac{\Delta
    \left(c_{\min}\right)^{3/2}}{2\left(c_{\max}\right)^{1/2}}.
  \end{equation}
  In particular, if $\overline{\bm{s}}_1\neq \overline{\bm{v}}^\star$ then the bound above is strict, so $\overline{\bm{v}}^\star$ is the unique maximizer of $L$ on $\overline{\mathcal{S}}$.

  \emph{Step 2: For fixed investment $\bm{c}^\top\bm{s}_1=\invst$ small enough, the unique optimal unit-investment design is $\overline{\bm{v}}^\star$.}
  By \linkedCref{prop:small-investment-expansion-appendix}, there exist $\overline{\invst}>0$ and $0<C_1<\infty$ such that for all $\invst\in(0,\overline{\invst}]$ and $\overline{\bm{s}}_1\in\overline{\mathcal{S}}$,
  \[
    V_Y(\invst\overline{\bm{s}}_1)
    =\alpha_1
    \E[Y]+\sqrt{\invst}\,\varphi(z_{1-\alpha_1})\,L(\overline{\bm{s}}_1)+R(\invst,\overline{\bm{s}}_1),
  \]
  with the $O(\invst)$-Lipschitz remainder bound
  \[
    |R(\invst,\overline{\bm{s}}_1)-R(\invst,\overline{\bm{s}}_1')|
    \le C_1\,\invst\,\|\overline{\bm{s}}_1-\overline{\bm{s}}_1'\|_1
    \qquad\forall\overline{\bm{s}}_1,\overline{\bm{s}}_1'\in\overline{\mathcal{S}}.
  \]
  Therefore, for any $\overline{\bm{s}}_1\in\overline{\mathcal{S}}$ and any $\invst\in(0,\overline{\invst}]$,
  \begin{align*}
    V_Y(\invst\overline{\bm{v}}^\star)-V_Y(\invst\overline{\bm{s}}_1)
    &=
    \sqrt{\invst}\,\varphi(z_{1-\alpha_1})\Big(L(\overline{\bm{v}}^\star)-L(\overline{\bm{s}}_1)\Big)
    +\Big(R(\invst,\overline{\bm{v}}^\star)-R(\invst,\overline{\bm{s}}_1)\Big)\\
    &\ge
    \sqrt{\invst}\,\varphi(z_{1-\alpha_1})\Big(L(\overline{\bm{v}}^\star)-L(\overline{\bm{s}}_1)\Big)
    -C_1\,\invst\,\|\overline{\bm{s}}_1-\overline{\bm{v}}^\star\|_1.
  \end{align*}
  Combining this with \eqref{eq:generic-leading-gap} gives
  \[
    V_Y(\invst\overline{\bm{v}}^\star)-V_Y(\invst\overline{\bm{s}}_1)
    \ge
    \Big(\sqrt{\invst}\,\varphi(z_{1-\alpha_1})\,q-C_1
    \invst\Big)\,\|\overline{\bm{s}}_1-\overline{\bm{v}}^\star\|_1.
  \]
  Let
  \[
    \invst_+ \defeq
    \min\!\left\{\overline{\invst},\ \Big(\frac{\varphi(z_{1-\alpha_1})q}{2C_1}\Big)^2\right\}.
  \]
  Then for all $\invst\in(0,\invst_+]$,
  \[
    \sqrt{\invst} \le \sqrt{\invst_+}
    \le \frac{\varphi(z_{1-\alpha_1})q}{2C_1}
  \]
  and the coefficient in parentheses satisfies
  \[
    \sqrt{\invst}\,\varphi(z_{1-\alpha_1})q-C_1 \invst
    \ge
    \sqrt{\invst}\,\varphi(z_{1-\alpha_1})q-C_1 \sqrt{\invst}
    \frac{\varphi(z_{1-\alpha_1})q}{2C_1}
    =
    \tfrac{1}{2}\sqrt{\invst}\,\varphi(z_{1-\alpha_1})q>0,
  \]
  so $V_Y(\invst\overline{\bm{v}}^\star)>V_Y(\invst\overline{\bm{s}}_1)$ for every $\overline{\bm{s}}_1\neq \overline{\bm{v}}^\star$ in $\overline{\mathcal{S}}$.
  Hence $\overline{\bm{v}}^\star$ is the unique maximizer of $\overline{\bm{s}}_1\mapsto V_Y(\invst\overline{\bm{s}}_1)$ on $\overline{\mathcal{S}}$ for all $\invst\in(0,\invst_+]$.

  \emph{Step 3: The optimal spending decision for sufficiently small budgets.}
  Let $\invst_+>0$ be as in \emph{Step 2}.
  Recall that $L(\overline{\bm{v}}^\star)=\pi_{d^\star}^Y>0$, and write
  \[
    A \;\defeq\; \varphi(z_{1-\alpha_1})\,L(\overline{\bm{v}}^\star)
    \;=\;\varphi(z_{1-\alpha_1})\,\pi_{d^\star}^Y \;>\;0.
  \]
  By \linkedCref{prop:small-investment-expansion-appendix}, the remainder is Lipschitz in the investment: for all $\invst,\invst'\in(0,\overline{\invst}]$,
  \[
    \big|R(\invst,\overline{\bm{v}}^\star)-R(\invst',\overline{\bm{v}}^\star)\big|
    \le C_0\,|\invst-\invst'|.
  \]
  Define
  \[
    \overline{b} \defeq
    \min\!\left\{\,\invst_+,\ \overline{\invst},\ \left(\frac{A}{2C_0}\right)^2\right\}.
  \]
  The map $\invst\mapsto V_Y(\invst\,\overline{\bm{v}}^\star)$ is strictly increasing on $(0,\overline{b}]$.
  To see this, fix any $0<\invst<\invst'\le \overline{b}$.
  Directly applying \linkedCref{prop:small-investment-expansion-appendix} with $\overline{\bm{s}}_1=\overline{\bm{v}}^\star$ gives
  \[
    V_Y(\invst'\overline{\bm{v}}^\star)-V_Y(\invst\overline{\bm{v}}^\star)
    =
    A(\sqrt{\invst'}-\sqrt{\invst})
    +\Big(R(\invst',\overline{\bm{v}}^\star)-R(\invst,\overline{\bm{v}}^\star)\Big).
  \]
  By the Lipschitz-in-$\invst$ remainder bound,
  \begin{align*}
    V_Y(\invst'\overline{\bm{v}}^\star)-V_Y(\invst\overline{\bm{v}}^\star)
    &\ge
    A(\sqrt{\invst'}-\sqrt{\invst}) - C_0(\invst'-\invst)
    = (\invst'-\invst)\left(\frac{A}{\sqrt{\invst'}+\sqrt{\invst}}-C_0\right).
  \end{align*}
  Since $\invst<\invst'$ we have $\sqrt{\invst'}+\sqrt{\invst}<2\sqrt{\invst'}$, hence
  \[
    \frac{A}{\sqrt{\invst'}+\sqrt{\invst}}
    >
    \frac{A}{2\sqrt{\invst'}}
    \ge
    \frac{A}{2\sqrt{\overline{b}}}
    \ge C_0,
  \]
  where the last inequality uses $\overline{b}\le (A/(2C_0))^2$.
  Therefore $V_Y(\invst'\overline{\bm{v}}^\star)>V_Y(\invst\overline{\bm{v}}^\star)$.

  Now fix any budget $b\in(0,\overline{b}]$.
  By \emph{Step 2}, for every fixed investment $\invst\in(0,b]$ the unique maximizer of $\overline{\bm{s}}_1\mapsto V_Y(\invst\overline{\bm{s}}_1)$ over $\overline{\mathcal{S}}$ is $\overline{\bm{v}}^\star$.
  Since $\invst\mapsto V_Y(\invst\overline{\bm{v}}^\star)$ is strictly increasing on $(0,b]$, the unique maximizer over $\invst\in(0,b]$ is $\invst=b$.
  Hence the unique optimizer is
  \[
    \bm{s}_1^\star(b)=b\,\overline{\bm{v}}^\star
    =
    \frac{b}{c_{d^\star}}\,\bm{e}_{d^\star}.
  \]
\end{proof}

\subsection{Named objectives under a general prior}
\label{sec:named-first-stage-results-appendix}

\begin{proposition}[Finite fourth moments imply the generic small-investment assumptions for the two-stage impact objective]
  \label{prop:fourth-moment-small-investment-two-impact-appendix}
  Suppose the prior distribution of $\bm\tau$ has finite fourth moments.
  Then the downstream payoff
  \[
    Y_{\twoimpact}=\Test_2\cdot \ATE_3
  \]
  is non-adaptive and satisfies $\E[|Y_{\twoimpact}|]<\infty$ and $\E[|Y_{\twoimpact}|\|\bm\tau\|_2^3]<\infty$.
  Consequently, the non-adaptivity and moment assumptions in \linkedCref{prop:small-investment-expansion-appendix} and \linkedCref{thm:optimal-small-budget-design-general-appendix} hold with $Y=Y_{\twoimpact}$.
\end{proposition}

\begin{proof}
  $Y_{\twoimpact}=\Test_2\cdot \ATE_3$ is non-adaptive because $\Test_2$ and $\ATE_3$ are independent of the pilot-stage sampling noise conditional on $\bm\tau$ by the probability model described in \Cref{sec:Model}.

  The first integrability condition is satisfied because
  \[
    |Y_{\twoimpact}|
    \le
    |\ATE_3|
    =
    \frac{|\bm s_3^\top\bm\tau|}{n_3}
    \le
    \frac{\|\bm s_3\|_2}{n_3}\|\bm\tau\|_2,
  \]
  and the second is satisfied because
  \[
    |Y_{\twoimpact}|\|\bm\tau\|_2^3
    \le
    |\ATE_3|\|\bm\tau\|_2^3
    =
    \frac{|\bm s_3^\top\bm\tau|}{n_3}\|\bm\tau\|_2^3
    \le
    \frac{\|\bm s_3\|_2}{n_3}\|\bm\tau\|_2^4,
  \]
  where we use $0\le\Test_2\le1$, the definition
  $\ATE_3=\bm s_3^\top\bm\tau/n_3$, and Cauchy--Schwarz.
  Since $D<\infty$,
  \[
    \|\bm\tau\|_2^4
    =
    \left(\sum_{d=1}^D \tau_d^2\right)^2
    \le
    D\sum_{d=1}^D \tau_d^4,
  \]
  so finite fourth moments of $\bm\tau$ imply $\E[\|\bm\tau\|_2^4]<\infty$, and hence also $\E[\|\bm\tau\|_2]<\infty$.
  Therefore the right-hand sides are integrable, which proves the claim.
\end{proof}

Applying \linkedCref{prop:small-investment-expansion-appendix} and \linkedCref{thm:optimal-small-budget-design-general-appendix} with $Y=Y_{\twoimpact}=\Test_2 \cdot \ATE_3$ nearly recovers the small-budget results for the paper's main objective, and applying them with $Y=1, \Test_2$, and $\ATE_3$ gives the corresponding results for the alternate objectives.
The only missing component is that \linkedCref{thm:optimal-small-budget-design-general-appendix} requires the maximizing type to have a positive index.
We now show that this positivity condition holds automatically for the two impact objectives $V_{\singleimpact}$ and $V_{\twoimpact}$.

\begin{proposition}[Automatic positivity for the impact objectives]
  \label{prop:positive-impact-objectives-appendix}
  Suppose the prior distribution of $\bm\tau$ has finite fourth moments.
  If there exists a unique type maximizing the small-budget index for $V_{\singleimpact}$, then there exists at least one type $d$ with $\theta_d^{\singleimpact}>0$.
  If there exists a unique type maximizing the small-budget index for $V_{\twoimpact}$, then there exists at least one type $d$ with $\theta_d = \theta^{\twoimpact}_d>0$.
  Consequently, the positivity condition in \linkedCref{thm:optimal-small-budget-design-general-appendix} is automatic for either impact objective whenever its small-budget index has a unique maximizer.

\end{proposition}

\begin{proof}
  To begin, note that if there exists a unique type $d^\star$ maximizing the small-budget index for a payoff $Y$, then
  \[
    \pi_d^Y = \frac{\theta_d^Y}{\sqrt{c_d}} = \frac{\E[\tau_d \cdot Y]}{\sqrt{c_d}}
  \]
  is not constant, and thus $Y$ cannot be almost surely zero.
  Also, the assumption that the prior has finite fourth moments implies that $\E[\ATE_3^2]<\infty$.

  For the single-impact objective, $Y_{\singleimpact}=\ATE_3$, so a unique small-budget index maximizer implies that $\ATE_3$ is not almost surely zero.
  Taking the inner product with $\bm s_3$ gives
  \[
    \bm s_3^\top\bm\theta^{\singleimpact}
    =
    \sum_{d=1}^D s_{3d}\,\E[\tau_d \ATE_3]
    =
    \E[(\bm s_3^\top\bm\tau)\,\ATE_3]
    =
    n_3\,\E[\ATE_3^2].
  \]
  Since $\ATE_3$ is not almost surely zero, we have $\E[\ATE_3^2]>0$, so $\bm s_3^\top\bm\theta^{\singleimpact}>0$.
  Therefore at least one coordinate of $\bm\theta^{\singleimpact}$ is strictly positive.

  For the two-impact objective, $Y_{\twoimpact}=\Test_2\cdot \ATE_3$, so a unique small-budget index maximizer implies that $\Test_2\cdot \ATE_3$ is not almost surely zero, and therefore $\ATE_3$ is not almost surely zero.
  Similarly,
  \[
    \bm s_3^\top\bm\theta^{\twoimpact}
    =
    \sum_{d=1}^D s_{3d}\,\E[\tau_d\,\cdot\,\Test_2\,\cdot\, \ATE_3]
    =
    \E[(\bm s_3^\top\bm\tau)\,\cdot\,\Test_2\,\cdot\, \ATE_3]
    =
    n_3\,\E[\Test_2\,\cdot\, \ATE_3^2].
  \]
  Conditional on $\bm\tau$, $\Test_2$ is Bernoulli with success probability $p_2(\ATE_2)\in(0,1)$.
  Hence
  \[
    \E[\Test_2\,\cdot\, \ATE_3^2]
    =
    \E[p_2(\ATE_2)\,\ATE_3^2]
    >
    0,
  \]
  because $p_2(\ATE_2)>0$ and $\ATE_3$ not almost surely zero implies $\P(\ATE_3\neq0)>0$.
  Thus $\bm s_3^\top\bm\theta^{\twoimpact}>0$, so at least one coordinate of $\bm\theta^{\twoimpact}$ is strictly positive.
\end{proof}

Combining \linkedCref{thm:optimal-small-budget-design-general-appendix} with \Cref{prop:fourth-moment-small-investment-two-impact-appendix} and \Cref{prop:positive-impact-objectives-appendix} gives \Cref{cor:optimal-small-budget-pilot}.

The small-budget indices can be written in the following form for general priors.
In \Cref{sec:elliptical-decompositions-appendix} we show that the two-stage objectives can be further simplified under an elliptical prior.

\begin{proposition}[Representations for the named coefficients under a general prior]
  \label{prop:first-stage-coefficient-formulas-appendix}
  Assume the prior on $\bm\tau$ has finite second moments.
  Then for every type $d$,
  \begin{align*}
    \theta_d^{\singlesuccess}
    &=
    \E[\tau_d] \\
    \theta_d^{\twosuccess}
    &=
    \E[\tau_d\,p_2(\ATE_2)],\\
    \theta_d^{\singleimpact}
    &=
    \E[\tau_d \ATE_3]
    =
    \E[\ATE_3]\,\E[\tau_d]
    +
    \cov(\tau_d,\ATE_3),\\
    \theta_d
    =
    \theta_d^{\twoimpact}
    &=
    \E[\tau_d\,p_2(\ATE_2)\,\ATE_3].
  \end{align*}

\end{proposition}

\begin{proof}
  The single-stage success formula is immediate from the definition of $\theta_d^{\singlesuccess}$.

  Since $p_2(\ATE_2)=\P(\Test_2=1\mid\bm\tau)\in(0,1)$ and $\ATE_3$ is measurable with respect to $\bm\tau$, by iterated expectations we have
  \[
    \theta_d^{\twosuccess}
    =
    \E[\tau_d\,\cdot\,\Test_2]
    =
    \E[\tau_d\,p_2(\ATE_2)],
    \qquad
    \theta_d
    =
    \E[\tau_d\,\cdot\,\Test_2\,\cdot\, \ATE_3]
    =
    \E[\tau_d\,p_2(\ATE_2)\,\ATE_3].
  \]
  For the single-stage impact objective,
  \[
    \theta_d^{\singleimpact}
    =
    \E[\tau_d \ATE_3]
    =
    \E[\tau_d]\E[\ATE_3]+\cov(\tau_d,\ATE_3)
  \]
  by the linearity of expectation and the definition of covariance.
\end{proof}
